%=============================================================================!
% 15.4  WHERE A RUN BEGINS                                                    !
%=============================================================================!
\section{Where a run begins}
\label{sec:cv-equilibration}

A run starts from positions and velocities that the ensemble would rarely
produce: a lattice for a liquid, a box of the wrong volume, a structure
just built. Its early part relaxes towards the ensemble, and an average
that includes it is shifted however long the run. This section finds how
large that shift is, and how to choose the start of the part to keep.

\subsection*{A cup of coffee}

Hot coffee left on a table cools at a rate proportional to how much
warmer than the room it is: $\dd T/\dd t = -(T - T_{\mathrm{room}})/\tau_{\mathrm r}$,
with $\tau_{\mathrm r}$ set by the cup and the air. This is the decay
equation of \cref{sec:ne-drag} for $T - T_{\mathrm{room}}$, with
$1/\tau_{\mathrm r}$ in place of $\gamma$, so
\begin{equation*}
   T(t) = T_{\mathrm{room}} + \bigl(T(0) - T_{\mathrm{room}}\bigr)
   \mathrm{e}^{-t/\tau_{\mathrm r}} .
\end{equation*}
A thermometer in the cup, read from the moment the coffee is poured and
averaged, gives a temperature above the room's. A run behaves the same
way: a quantity $A$ approaches its ensemble average with a decaying offset
$\delta A_0\,\mathrm{e}^{-t/\tau_{\mathrm r}}$, which is $\delta A_0$ at the
start, and on which the run's own
fluctuations ride. Averaged over a run of length $t_{\mathrm f}$ from its
start, the offset contributes
\begin{equation*}
   \frac{1}{t_{\mathrm f}}\int_0^{t_{\mathrm f}}\delta A_0\,
   \mathrm{e}^{-t/\tau_{\mathrm r}}\,\dd t
   = \frac{\delta A_0\,\tau_{\mathrm r}}{t_{\mathrm f}}\Bigl[-\mathrm{e}^{-t/
   \tau_{\mathrm r}}\Bigr]_0^{t_{\mathrm f}}
   = \frac{\delta A_0\,\tau_{\mathrm r}}{t_{\mathrm f}}\bigl(1 -
   \mathrm{e}^{-t_{\mathrm f}/\tau_{\mathrm r}}\bigr)
   \approx \frac{\delta A_0\,\tau_{\mathrm r}}{t_{\mathrm f}}
\end{equation*}
for a run much longer than $\tau_{\mathrm r}$. The shift falls as
$1/t_{\mathrm f}$ and the error bar as $1/\sqrt{t_{\mathrm f}}$, so a long
enough run would hide the shift inside its error bar, but only slowly.
Leaving out the first $t_0$ of the run cuts the offset's share to
$\delta A_0\,\tau_{\mathrm r}\mathrm{e}^{-t_0/\tau_{\mathrm r}}/(t_{\mathrm f}
- t_0)$ (\cref{ex:cv-offset}), at the cost of the samples discarded.

\subsection*{The start that leaves the most}

Discarding too little leaves a shift; discarding too much throws away
samples. Chodera found a way to weigh the two from the run
alone\cite{Chodera2016}. To the autocorrelation function a decaying offset
looks like a very slow correlation, and it inflates $g$; so for each
candidate start $t_0$, the effective number of samples left after it,
\begin{equation*}
   n_{\mathrm{eff}}(t_0) = \frac{n - t_0/\dt}{g(t_0)} ,
\end{equation*}
with $t_0/\dt$ the number of samples discarded and $g(t_0)$ the
statistical inefficiency of the samples from $t_0$ on,
rises while the discard removes more of the offset's inflation of $g$
than it costs in samples, and falls once only samples are being removed.
The start kept is the one that maximises it:

\begin{code}
def detect_equilibration(
    series: ArrayLike, nskip: int = 1
) -> tuple[int, float, float]:
    a = np.asarray(series, dtype=float)
    n = len(a)
    best = (0, 1.0, 0.0)
    for t0 in range(0, n - 1, nskip):
        rest = a[t0:]
        if np.all(rest == rest[0]):
            g = float(n - t0)
        else:
            g = statistical_inefficiency(rest)
        if (n - t0) / g > best[2]:
            best = (t0, g, (n - t0) / g)
    return best
\end{code}

\noindent The listing counts \texttt{t0} in samples. It finds the same
start as the function of the same name in the \texttt{timeseries} module
of pymbar with \texttt{fast=False} (\texttt{test\_stats.py}); pymbar's
default estimates $g$ more roughly and can place the start a few samples
later, 550 against 534 for the series below.

\Cref{fig:cv-equilibration}(a) adds the offset $3\,\mathrm{e}^{-t/\SI{2}{\pico
\second}}$, three standard deviations at the start, to \SI{50}{\pico
\second} of the series of \cref{sec:cv-correlation}, and (b) shows
$n_{\mathrm{eff}}(t_0)$ rising to 158 at \SI{5.34}{\pico\second}, 2.7 decay
times, and falling after. From the start the mean is $0.160 \pm 0.164$,
of which the offset contributes $\delta A_0\,\tau_{\mathrm r}/t_{\mathrm f} = 0.120$; after
the discard it is $0.001 \pm 0.076$, the offset contributing 0.0093, and
the error bar has halved. One series is one draw: over 200 such series the
start falls between 0.1 and 16.5 decay times, with the median, the value
with half the series on either side, at 1.4, and
the offset left in the mean is at most 1.35 error bars, with the median
0.34. One error bar then holds the true mean in 0.58 of the series, where
it should hold it in 0.68; without the discard, in 0.54. The method trades
a shift of about a third of an error bar for an error bar half as large,
and leaves part of a slow offset in place, which the test for drift of
\cref{sec:cv-drift} looks for.

\begin{figure}[htb]
\centering
\includegraphics{ch15_convergence/equilibration}
\caption{Where a run begins. (a) \SI{50}{\pico\second} of the series of
\cref{eq:cv-ou} with $c = 0.95$, with the offset $3\,\mathrm{e}^{-t/\SI{2}{\pico
\second}}$ added, the offset itself dashed, and the start found (dotted).
(b) The effective number of samples after each candidate start,
maximised at the dotted line. (c) The first \SI{30}{\pico\second} of eight
runs of the liquid melting from the face-centred
cubic lattice at \SI{135}{\kelvin}: the potential energy per atom (grey),
their mean (teal) and the mean of the \SI{2}{\nano\second} run (dashed),
with the start found in each run from its own $U$ (ticks).}
\label{fig:cv-equilibration}
\end{figure}

\subsection*{A liquid melting}

\Cref{fig:cv-equilibration}(c) starts eight runs of the liquid from the
face-centred cubic lattice at the liquid's density, with velocities drawn
at \SI{135}{\kelvin} from eight seeds, under CSVR for \SI{100}{\pico
\second}. The lattice has $U/N = \SI{-64.331}{\milli\electronvolt}$; the
energy climbs from it as the lattice melts, and the mean of the eight runs
comes
within \SI{1}{\milli\electronvolt} of the \SI{2}{\nano\second} run's
$(-50.092 \pm 0.018)$\,\si{\milli\electronvolt} after \SI{2.89}{\pico
\second}. In four of the runs the start found from $U$ is 2.8 to
\SI{3.2}{\pico\second}, at the end of the melting. In the other four it is
19 to \SI{27}{\pico\second}: under CSVR the energy wanders over
picoseconds (\cref{sec:cv-correlation}), and a run of \SI{100}{\pico
\second} cannot always tell such a wander from the tail of a transient.
Discarding more than needed costs samples and, on average, does not shift
the mean.
Averaged from their first steps, the eight means of $U/N$ lie at $-50.169$
on average, the melting's share pulling them down; after the discards
they average \SI{-50.063}{\milli\electronvolt}, close to the
\SI{2}{\nano\second} value, and scatter by \SI{0.111}{\milli\electronvolt}
against error bars of 0.046 to \SI{0.085}{\milli\electronvolt}, 0.067 on
average. The error bars are low by about two fifths, though a spread from
eight runs is itself uncertain by a quarter (\cref{ex:cv-error-error}). A
run of \SI{100}{\pico\second} holds about 68 independent samples of $U$
(\cref{sec:cv-blocks}), and the start chosen makes $(n - t_0/\dt)/g(t_0)$
largest, which favours a start after which $g$ happens to come out small:
the eight runs find $g$ of 45 to 138 after their starts, against the 146
of the \SI{2}{\nano\second} run, and their error bars are small with it.

The same holds for any quantity with few independent samples.
Chapter~14's three runs of \SI{300}{\pico\second} at constant pressure
under stochastic cell rescaling hold 35, 73 and 109 independent samples of
the volume. In the second and third the start found is 1.0 and
\SI{0.0}{\pico\second}, well within the first \SI{10}{\pico\second} that
\cref{sec:pr-trajectory} left out of its averages. In the first,
$n_{\mathrm{eff}}(t_0)$ falls from 35 at the start, as a settled run's
should, to 24 at \SI{100}{\pico\second} and 9 at \SI{195}{\pico\second};
beyond, the $g$ estimated from the last few tens of picoseconds is so
uncertain that it can come out far too small, and the largest value, 55,
falls at \SI{259}{\pico\second}, a start with no meaning. The criterion of \cref{sec:cv-protocol}
therefore asks that the start lie in the first half of the run. Chapter~11
prepared its water in rounds of \SI{1}{\pico\second} and left the end of
that slow change unmeasured (\cref{sec:cs-water}); its healthy run of
\SI{10}{\pico\second} at fixed energy has its start at \SI{0.30}{\pico
\second}, after which the box's kinetic energy is $K = (4.993 \pm
0.022)$\,\si{\electronvolt} and the potential energy $(-0.40265 \pm
0.00034)$\,\si{\electronvolt} per molecule, against 4.990 and $-0.40259$
over the whole run: the preparation had all but finished, and
the averages that \cref{sec:cs-trajectory} quoted stand within their error
bars.

\begin{keyidea}
A run's early part relaxes from its start, and shifts the mean of the
whole run by about $\delta A_0\,\tau_{\mathrm r}/t_{\mathrm f}$. The start to keep is the
one that leaves the largest effective number of samples, $(n - t_0/\dt)/
g(t_0)$. The method trades a shift of a fraction of an error bar for
precision and needs a
run with many independent samples: when the start it finds falls in the
second half of the run, the run is too short to locate it.
\end{keyidea}

\notebook{15}{add a transient of your choosing and watch where the start
is found}

\begin{selfcheck}
A run of \SI{100}{\pico\second} has an offset of \SI{2}{\kelvin} in its
temperature at the start, decaying with the time $\tau_{\mathrm r} =
\SI{1}{\pico\second}$. By how much
does it shift the mean temperature of the whole run?
\end{selfcheck}
